\documentclass{article}


% Title Page
\title{\Huge \textbf{Simulazioni \textit{n}-body: vari approcci paralleli e approssimati}}
\author{Relazione del progetto di GPGPU - Damiano Trovato}
\date{}

\usepackage{graphicx} % Per le immagini
\usepackage{amsmath} % Simboli matematici
\usepackage[hidelinks]{hyperref} % Collegamenti ipertestuali
\usepackage{amssymb} % Simboli matematici 2
\usepackage{tikz}
\usepackage{fancyvrb}
\usepackage{multicol}
\usepackage{geometry}
\usepackage{array}
\usepackage{tabularx} % nel preambolo
\usepackage{listings}
\usepackage{algorithm}
\usepackage{mathabx}
\usepackage{algpseudocode}

\renewcommand{\thelstnumber}{\arabic{lstnumber}:}
\lstdefinestyle{mystyle}{          
	numbers=left,                    
	numbersep=5pt             
}
\lstset{style=mystyle}

\makeatother
\geometry{
	a4paper,
	total={160mm,260mm},
	left=20mm,
	top=18mm,
}
% RIDENOMINAZIONI
\renewcommand{\contentsname}{Indice}
\renewcommand{\figurename}{Immagine}


\begin{document}

	\maketitle
	\begin{abstract}
	Con la seguente relazione si documenta il processo di realizzazione di una simulazione $n$-body tramite soluzioni parallele su GPU. Si trattano quindi aspetti degni di nota, quali approcci, metodi numerici, ottimizzazioni e analisi delle performance. Si offre quindi un'analisi progressiva di più soluzioni, partendo da una soluzione seriale\footnote{L'analisi delle soluzioni seriali sarà più sintetica, in quanto non riflette totalmente lo scope della materia per cui si realizza questo progetto.} su CPU, a una parallela GPU naive $\mathcal{O}(n^2)$, fino ad una soluzione seriale approssimata basata sull'algoritmo Barnes-Hut $\mathcal{O}(n\log n)$, e, definitivamente, la sua controparte parallela su GPU. 
	\end{abstract}
	
	\section{Dettagli generali}
	In questa prima parte della relazione, si trattano gli aspetti comuni a tutte le soluzioni trattate.
	
	\subsection{Aspetti matematici}
	Si offre una lista delle formule matematiche che rendono fattibili le simulazioni realizzate all'interno del progetto.
	
	\paragraph{Legge di gravitazione universale} La seguente formula detta l'attrazione che agisce tra due corpi rispettivamente di massa $m_1$ ed $m_2$, situate a distanza $R$. Questa grandezza è direttamente proporzionale al prodotto delle masse e inversamente proporzionale al quadrato della distanza, tutto moltiplicato per una costante (di cui verrà modificato il valore ai fini della simulazione).
	$$
	F = G \frac{m_1m_2}{R^2}
	$$
	
	\paragraph{Dalla forza all'accelerazione}
	
	$$
	F = ma \Rightarrow a = \frac{F}{m}
	$$
	
	\paragraph{Dall'accelerazione alla velocità} L'accelerazione è la derivata della velocità sul tempo, possiamo quindi ottenere
	
	$$
	a = \frac{dv}{dt} \Rightarrow adt = \frac{F}{m}dt = dv
	$$
	
	fissando un valore di $dt$, tramite la variabile \verb*|delta_time|, andremo a integrare rispetto al tempo, ottenendo la velocità.
	
	\paragraph{Dalla velocità allo spostamento}
	
	$$
	x(t) = x_0 + v dt
	$$
	
	dove $dt = \verb*|delta_time|$.
	
	\paragraph{Lavorare in tre dimensioni} Lavorando in tre dimensioni, il calcolo di ciascuna delle forze, velocità e spostamenti è sempre relativo alle singole componenti $(x,y,z)$. Nel calcolo delle forze, ad esempio, andiamo a moltiplicare la forza ottenuta per ciascuna componente del versore (ottenuto normalizzando il vettore calcolato). Da qui in poi, otteniamo la velocità per ciascuna delle componenti, e quindi lo spostamento nello spazio tridimensionale.
	
	
	\subsection{Metodo numerico di integrazione: leapfrog}
	Il metodo numerico utilizzato nella simulazione, è noto come metodo \textbf{leapfrog}, o metodo della cavallina. Il nome deriva dal modo in cui si alterna l'integrazione della velocità a quella del tempo.
	
	\begin{algorithm}
		\caption{Integrazione Leapfrog}
		\begin{algorithmic}[1]
		\State Calcola la forza
		\State Aggiorna la velocità per $\Delta t/2$
		\For{Numero di iterazioni}
			\State Aggiorna la posizione per $\Delta t$
			\State Calcola la forza
			\State Aggiorna la velocità per $\Delta t$
		\EndFor
		\end{algorithmic}
	\end{algorithm}
	
	Questo metodo numerico, sfasando di $\Delta t/2$ l'integrazione della velocità e della posizione, ottiene, nel contesto delle simulazioni $n$-body, risultati più accurati a parità di operazioni. L'accuratezza guadagnata, rispetto a un banale metodo di Eulero, si manifesta nell'energia del sistema, che viene conservata (con oscillazioni trascurabili).
	
	
	\section{Simulazione seriale naive}	
	Il primo approccio alla simulazione di sistemi a $n$ corpi che si è deciso di realizzare, è una soluzione seriale basata sull'algoritmo naive, con complessità computazionale $\mathcal{O}(n^2)$.
	
	\section{Simulazione parallela naive su GPU}
	
	\section{Simulazione approssimata Barnes-Hut seriale}
	
	
	\section{Simulazione approssimata Barnes-Hut su GPU}
	
\end{document}          
